\documentclass[slideColor]{prosper}
\usepackage{amsmath,amsfonts,amsthm}
%\usepackage{epsfig}
\include{Qcircuit}


\newcommand{\E}{\mathbf{E}}
%\newcommand{\R}{\mathbb{R}}

\newcommand{\cA}{\mathcal A}
\newcommand{\cB}{\mathcal B}
\newcommand{\cH}{\mathcal H}
\newcommand{\C}{\mathbb C}

\renewcommand{\>}{\rangle}
\newcommand{\<}{\langle}


%%%%%%%%%%%%%%%%%%%%%%%%%%%%%%%%%%%%%%%%%%%%%%%%%%%%%%%%%%%%

\title{Quantum Speed-Up of Estimating Partition Functions}
\author{Pawel M. Wocjan\\[1ex]School of Electrical Engineering and Computer Science\\ University of Central Florida\\Orlando}
\email{wocjan@eecs.ucf.edu}

\begin{document}

\maketitle

%%%%%%%%%%%%%%%%%%%%%%%%%%%%%%%%%%%%%%%%%%%%%%%%%%%%%%%%%%%%%%%%%%%%%%%%

\begin{slide}{Joint work with}
\begin{itemize}
\item Chen-Fu Chiang (University of Central Florida)

\bigskip
\item Anura Abeyesinghe (University of Central Florida)

\bigskip
\item Daniel Nagaj (Slovak Academy of Sciences)

\bigskip\bigskip
\item this work has been supported by the National Science Foundation
grants CCF-0726771 and CCF-0746600
\end{itemize}
\end{slide}

%%%%%%%%%%%%%%%%%%%%%%%%%%%%%%%%%%%%%%%%%%%%%%%%%%%%%%%%%%%%%%%%%%%%%%%%%

\begin{slide}{Related Articles}
\begin{itemize}
\item Szegedy, {\em Spectra of Quantized Walks and a $\sqrt{\delta\epsilon}$ rule}, 2004

\medskip
\item Magniez, Nayak, Roland, Santha, {\em Search via Quantum Walk}, 2006

\medskip
\item Santha, {\em Quantum Walk Based Search Algorithms}, 2008
(overview article)

\medskip
\item Somma, Boixo, Barnum, Knill, {\em Quantum Simulations of Classical Annealing Processes}, 2008

\medskip
\item W., Abeyesinghe, {\em Speed-Up via Quantum Sampling}, 2008
\end{itemize}

\end{slide}

%%%%%%%%%%%%%%%%%%%%%%%%%%%%%%%%%%%%%%%%%%%%%%%%%%%%%%%%%%%%%%%%%%%%%%%%%

\begin{slide}{Estimating Partition Functions}
\begin{itemize}
\item let $\Omega$ be a set whose elements $x$ correspond the states of some physical system

\bigskip
\item let $E : \Omega\rightarrow \mathbb{R}$ denote the energy function, assigning each state $x$ its energy $E(x)$

\bigskip
\item given the desired (inverse) temperature $\beta$, the task is to estimate the partition function $Z(\beta)$

\[
Z(\beta) = \sum_{x\in\Omega} e^{-\beta E(x)}
\]
\end{itemize}

\end{slide}

%%%%%%%%%%%%%%%%%%%%%%%%%%%%%%%%%%%%%%%%%%%%%%%%%%%%%%%%%%%%%%%%%%%%%%%%%

\begin{slide}{Computational Complexity}
\begin{itemize}

\bigskip
\item the problem of estimating $Z(\beta)$ with \textcolor{red}{high degree of accuracy} for \textcolor{blue}{general} $E : \Omega\rightarrow \mathbb{R}$ and \textcolor{green}{high} $\beta$ is \#P-hard

\bigskip
\item $\Longrightarrow$ It is unlikely that there are efficient (classical and quantum) algorithms for this task 
\end{itemize}

\end{slide}

%%%%%%%%%%%%%%%%%%%%%%%%%%%%%%%%%%%%%%%%%%%%%%%%%%%%%%%%%%%%%%%%%%%%%%%%%


\begin{slide}{FPRAS}
\begin{itemize}
\item we consider fully polynomial randomized approximation schemes (FPRAS)

\medskip
\item a FPRAS

\begin{itemize}
\item outputs a random number $\hat{Z}$ satisfying
\[
\Pr\Big( (1-\epsilon)Z(\beta) \le \hat{Z} \le (1+\epsilon) Z(\beta) \Big) \ge 3/4
\]
where $\epsilon\in (0,1)$ determines the desired accuracy

\medskip
\item runs in time polynomial in $\log(|\Omega|)$ and $1/\epsilon$
\end{itemize}
\end{itemize}
\end{slide}

%%%%%%%%%%%%%%%%%%%%%%%%%%%%%%%%%%%%%%%%%%%%%%%%%%%%%%%%%%%%%%%%%%%%%%%%%

\begin{slide}{Simulated Annealing}
\begin{itemize}
\item choose a cooling schedule $\beta_0 < \beta_1 < \cdots < \beta_{\ell}$ with $\beta_0=0$ and $\beta_{\ell}=\beta$

\medskip
\item express the desired quantity as a telescoping product
\[
Z(\beta) = \frac{Z(\beta_{\ell})}{Z(\beta_{\ell-1})} \cdot \frac{Z(\beta_{\ell-1})}{Z(\beta_{\ell-2})} \cdots \frac{Z(\beta_1)}{Z(\beta_0)} \cdot Z(\beta_0) 
\]

\medskip
\item observe that $Z(\beta_0)$ is trivial since $Z(\beta_0)=|\Omega|$

\medskip
\item $\Rightarrow$ estimate the ratios $\alpha_i:=Z(\beta_{i+1})/Z(\beta_i)$ 
\end{itemize}
\end{slide}

%%%%%%%%%%%%%%%%%%%%%%%%%%%%%%%%%%%%%%%%%%%%%%%%%%%%%%%%%%%%%%%%%%%%%%%%%

\begin{slide}{Estimation via Boltzmann Sampling I}
\begin{itemize}
\item denote by $\pi_i=\big(\pi_i(x) : x\in\Omega\big)$ the Boltzmann distribution at inverse temperature $\beta_i$
\[
\pi_i(x) = \frac{e^{-\beta_i E(x)}}{Z(\beta_i)}
\]

\item assume we can sample from $\pi_i$

\medskip
\item assume we can find a short cooling schedule so that each ratio $\alpha_i:=Z(\beta_{i+1})/Z(\beta_i)$ is bounded from below by a constant, say $1/2$

\medskip
\item $\Rightarrow$ it suffices to take a small number of samples $\sim \pi_i$ to estimate $\alpha_i$

\end{itemize}
\end{slide}

%%%%%%%%%%%%%%%%%%%%%%%%%%%%%%%%%%%%%%%%%%%%%%%%%%%%%%%%%%%%%%%%%%%%%%%%%

\begin{slide}{Estimation via Boltzmann Sampling II}
\begin{itemize}
\item let $X_i \sim \pi_i$

\medskip
\item let $\Delta\beta_i = \beta_{i+1} - \beta_i$ 

\medskip
\item define a new random variable 
\[
Y_i = e^{-\Delta\beta_i E(X_i)}
\]
 
\item the expected value $\E(Y_i)$ is equal to
\[
\sum_{x \in \Omega}
\pi_i(x) 
e^{-\Delta\beta_i E(x)}  =
\sum_{x \in \Omega}
\frac{e^{-\beta_i E(x)}}{Z(\beta_i)} 
e^{(-\beta_{i+1} + \beta_i) E(x)}  
= 
\alpha_i
\]
\item $\Rightarrow$ $Y_i$ is an unbiased estimator for $\alpha_i$
\end{itemize}
\end{slide}

%%%%%%%%%%%%%%%%%%%%%%%%%%%%%%%%%%%%%%%%%%%%%%%%%%%%%%%%%%%%%%%%%%%%%%%%%

\begin{slide}{Estimation via Boltzmann Sampling III}
\begin{itemize}
\item draw $O(\ell/\epsilon^2)$ samples of $X_i$ and compute the mean $\bar{Y}_i$

\medskip
\item $\Rightarrow$ the random variable $\bar{Y}=\bar{Y}_0 \bar{Y}_1 \cdots \bar{Y}_{\ell}$ satisfies
\[
\Pr\Big(
(1-\epsilon) \alpha < \bar{Y} < (1+\epsilon) \alpha
\Big)
\ge 7/8
\]
where $\alpha=\alpha_0\alpha_1\cdots\alpha_\ell$

\medskip
\item $\Rightarrow$ the total number of samples is
\[
O\Big(\ell^2 / {\red \epsilon^2} \Big)
\] 
\end{itemize}
\end{slide}

%%%%%%%%%%%%%%%%%%%%%%%%%%%%%%%%%%%%%%%%%%%%%%%%%%%%%%%%%%%%%%%%%%%%%%%%%

\begin{slide}{Sampling with Markov Chains}
\begin{itemize}
\item in general, we are not able to sample directly from $\pi_i$ 

\medskip
\item for some $E : \Omega \rightarrow \mathbb{R}$, we can construct a Markov chain $P_i$ such that 

\begin{itemize}
\medskip
\item its stationary distribution is equal to $\pi_i$ 

\medskip
\item its spectral gap $\delta$ is large
\end{itemize}

\medskip
\item $\Rightarrow$ we simulate  
\[
\tilde{O}(1/ {\red \delta})
\] 
steps of $P_i$ to obtain one sample from $\tilde{\pi}_i$, which is sufficiently close to $\pi_i$
\end{itemize}
\end{slide}

%%%%%%%%%%%%%%%%%%%%%%%%%%%%%%%%%%%%%%%%%%%%%%%%%%%%%%%%%%%%%%%%%%%%%%%%%

\begin{slide}{Quantum Speed-Up}
\begin{itemize}
\item classical complexity
\[
\tilde{O}\Big(
\ell^2 /\big(\delta \epsilon^2\big) 
\Big)
\]

\item quantum complexity
\[
\tilde{O}\Big(
\ell^2/\big(\sqrt{\delta} \epsilon\big)
\Big)
\]

\item $1/\delta \rightarrow 1\sqrt{\delta}$ is due to the quadratic relation between the spectral gaps of $P_i$ and the phase gaps of the corresponding quantum walks $W(P_i)$ 

\medskip
\item $\epsilon^2 \rightarrow \epsilon$ is due to quantum estimation of expected values of observables
\end{itemize}
\end{slide}

%%%%%%%%%%%%%%%%%%%%%%%%%%%%%%%%%%%%%%%%%%%%%%%%%%%%%%%%%%%%%%%%%%%%%%%%%

\begin{slide}{Classical Walk}
\begin{itemize}
\item $P_i$ is a time-reversible Markov chain 
\[
P_i=\Big(p_i(x,y)\Big)_{x,y\in\Omega}
\]
with stationary distribution 
\[
\pi_i=\Big(\pi_i(x)\Big)_{x\in\Omega}
\]

\item $\pi_i$ is the unique left eigenvector of $P_i$ with eigenvalue $1$
\[
\pi_i P_i = \pi_i
\]

\end{itemize}
\end{slide}

%%%%%%%%%%%%%%%%%%%%%%%%%%%%%%%%%%%%%%%%%%%%%%%%%%%%%%%%%%%%%%%%%%%%%%%%%

\begin{slide}{Spectral Gap}

\begin{itemize}

\item let $1=\lambda_0 > \lambda_1 > \ldots > \lambda_{|\Omega|-1}>0$ be the spectrum of $P_i$

\bigskip
\item the spectral gap 
\[
\delta := 1 - \lambda_1
\]

\medskip
determines how fast $\pi_i$ is approached from any initial state
\end{itemize}
\end{slide}

%%%%%%%%%%%%%%%%%%%%%%%%%%%%%%%%%%%%%%%%%%%%%%%%%%%%%%%%%%%%%%%%%%%%%%%%%

\begin{slide}{Quantum Update}

\begin{itemize}
\item a quantum update for $P_i$ is any unitary $U_i$ acting on $\C^{|\Omega|}\otimes\C^{|\Omega|}$ with
\begin{equation*}
U_i \Big( |x\> \otimes |0\> \Big) = 
|x\> \otimes 
\sum_{y\in\Omega} \sqrt{p_i(x,y)} |y\> 
\end{equation*}
for some fixed state $0\in\Omega$

\bigskip
\item classical complexity: the number of times we apply $P_i$

\bigskip
\item quantum complexity: the number of times we apply $U_i$ or $U_i^\dagger$
\end{itemize}
\end{slide}

%%%%%%%%%%%%%%%%%%%%%%%%%%%%%%%%%%%%%%%%%%%%%%%%%%%%%%%%%%%%%%%%%%%%%%%%%

\begin{slide}{Quantum Walk}
\begin{itemize}
\item the quantum walk $W(P_i)$ is realized by applying $U_i$ twice and $U_i^\dagger$ twice (and also some other simple unitaries)

\bigskip
\item the unique eigenvector of $W(P_i)$ with eigenvalue $1$ is
\[
|\pi_i\> = \sum_{x\in\Omega} \sqrt{\pi_i(x)} |x\>
\]

\item
$|\pi_i\>$ is a coherent version of the limiting distribution $\pi_i$
\end{itemize}
\end{slide}

%%%%%%%%%%%%%%%%%%%%%%%%%%%%%%%%%%%%%%%%%%%%%%%%%%%%%%%%%%%%%%%%%%%%%%%%%

\begin{slide}{Quadratic Relation}
\begin{itemize}
\item the phase gap $\Delta$ of $W(P_i)$ is
\[
\min\{ |\varphi| : \mbox{ $e^{i\varphi}$ is an eigenvalue of $W(P_i)$, $e^{i\varphi}\neq 1$} \} 
\]

\medskip
\item the quadratic relation between the phase and spectral gaps
\[
\Delta \ge \sqrt{\delta}
\]
is at the heart of quantum speed-ups of many search problems
\end{itemize}
\end{slide}

%%%%%%%%%%%%%%%%%%%%%%%%%%%%%%%%%%%%%%%%%%%%%%%%%%%%%%%%%%%%%%%%%%%%%%%%%

\begin{slide}{Structure of the Quantum Algorithm}

\vspace{1cm}
\hspace{6cm}
\Qcircuit @C=0.9em @R=0.9em {
\lstick{|\pi_0\>}                                                                                  & \gate{\tilde{\alpha}_0}      & \qw \\ 
\lstick{|\pi_0\> \rightarrow |\tilde{\pi}_{1}\>}                                                   & \gate{\tilde{\alpha}_1}      & \qw \\ 
                                                                                                   &                   &     \\
\lstick{|\pi_0\> \rightarrow |\tilde{\pi}_{1}\> \rightarrow \cdots \rightarrow |\tilde{\pi}_{\ell-1}\>} & \gate{\tilde{\alpha}_{\ell-1}} & \qw    
}
\end{slide}

%%%%%%%%%%%%%%%%%%%%%%%%%%%%%%%%%%%%%%%%%%%%%%%%%%%%%%%%%%%%%%%%%%%%%%%%%%
%
%\begin{slide}{Structure of the Quantum Algorithm II}
%\vspace{-0.25cm}
%\begin{itemize}
%\item to prepare $|\tilde{\pi}_i\>$, we apply operators in $\{W(P_k) : k=0,\ldots,i\}$
%\[
%\tilde{O}\Big(
%\frac{i}{\sqrt{\delta}}
%\Big)
%\]
%times
%
%\item to obtain $\tilde{\alpha}_i$, we apply $W(P_i)$
%\[
%\tilde{O}\Big(
%\frac{\ell}{\sqrt{\delta} \epsilon}
%\Big)
%\] 
%times to the state $|\tilde{\pi}_i\>$
%
%\item
%$\Rightarrow$ the overall complexity is
%\vspace{-0.5cm}
%\[
%{}\quad\quad\quad\quad\tilde{O}\Big(
%\frac{\ell^2}{\sqrt{\delta} \epsilon}
%\Big)
%\] 
%\end{itemize}
%\end{slide}
%
%%%%%%%%%%%%%%%%%%%%%%%%%%%%%%%%%%%%%%%%%%%%%%%%%%%%%%%%%%%%%%%%%%%%%%%%%%

\begin{slide}{Preparation of Quantum Samples I}
\begin{itemize}
\item the fact that the ratios $\alpha_i$ are bounded from below by $1/2$ implies 
\[
|\<\pi_i|\pi_{i+1}\>|^2 \ge \frac{1}{2}
\]

\medskip
\item $\Rightarrow$ we can drive $|\pi_i\>$ to $|\tilde{\pi}_{i+1}\>$ by applying $W(P_i)$ and $W(P_{i+1})$
\[
\tilde{O}\Big(\frac{1}{\sqrt{\delta}}\Big)
\] 
times

\medskip
\item this is based on Grover's $\frac{\pi}{3}$-fixed point search

\end{itemize}
\end{slide}

%%%%%%%%%%%%%%%%%%%%%%%%%%%%%%%%%%%%%%%%%%%%%%%%%%%%%%%%%%%%%%%%%%%%%%%%%

\begin{slide}{Preparation of Quantum Samples II}
\begin{itemize}
\item starting from $|\pi_0\>$ (uniform superposition), we can prepare any $|\tilde{\pi}_i\>$ by
\[
|\pi_0\> \rightarrow |\tilde{\pi}_1\> \rightarrow \cdots \rightarrow |\tilde{\pi}_i\> \rightarrow \cdots \rightarrow |\tilde{\pi}_{\ell-1}\>
\]

\medskip
\item $\Rightarrow$ it suffices to apply operators in $\{W(P_k) : k = 0,\ldots,i\}$
\[
\tilde{O}\Big(\frac{i}{\sqrt{\delta}}\Big)
\]
times

\medskip
\item $\Rightarrow$ the complexity of preparing $|\tilde{\pi}_0\>\otimes |\tilde{\pi}_1\> \otimes \cdots \otimes |\tilde{\pi}_{\ell-1}\>$ is
\[
\tilde{O}\Big(
\frac{\ell^2}{\sqrt{\delta}}
\Big)
\]
\end{itemize}
\end{slide}

%%%%%%%%%%%%%%%%%%%%%%%%%%%%%%%%%%%%%%%%%%%%%%%%%%%%%%%%%%%%%%%%%%%%%%%%%

\begin{slide}{Quantum Estimation of the Ratios I}
\begin{itemize}
\item let $A_i$ be the observable
\[
A_i = \sum_{x\in\Omega} e^{-\Delta\beta_i E(x)} |x\>\<x|
\]

\medskip
\item we have 
\[
\<\pi_i|A_i|\pi_i\> = \alpha_i
\]

\medskip
where 
\[
|\pi_i\> = \sum_{x\in\Omega} \sqrt{\pi_i(x)} |x\>
\]

\end{itemize}
\end{slide}

%%%%%%%%%%%%%%%%%%%%%%%%%%%%%%%%%%%%%%%%%%%%%%%%%%%%%%%%%%%%%%%%%%%%%%%%%

\begin{slide}{Quantum Estimation of the Ratios II}
\begin{itemize}
\item let 
\[
V_i := \sum_{x\in\Omega} |x\>\<x| \otimes
\left(
\begin{array}{cc}
 \sqrt{e^{-\Delta\beta_i E(x)}}   & \sqrt{1-e^{-\Delta\beta_i E(x)}} \\
-\sqrt{1-e^{-\Delta\beta_i E(x)}} & \sqrt{e^{-\Delta\beta_i E(x)}}
\end{array}
\right)
\]

\medskip
\item let 
\[
|\psi_i\> = V_i \Big( |\pi_i\> \otimes |0\> \Big) \quad \mbox{ and } \quad \Gamma := I \otimes |0\>\<0|
\]

\medskip
\item $\Rightarrow$ we have
\begin{eqnarray*}
\<\psi_i|\Gamma|\psi_i\> & = & \<\pi_i|A_i|\pi_i\> \\
& = & \alpha_i
\end{eqnarray*}
\end{itemize}
\end{slide}

%%%%%%%%%%%%%%%%%%%%%%%%%%%%%%%%%%%%%%%%%%%%%%%%%%%%%%%%%%%%%%%%%%%%%%%%%

\begin{slide}{Quantum Estimation of the Ratios III}
\begin{itemize}

\item we obtain a random variable $Q_i$ with
\[
\hspace{-0.5cm}
\Pr\left( \Big(1-\frac{\epsilon}{2\ell}\Big)\alpha_i \le Q_i \le \Big(1+\frac{\epsilon}{2\ell}\Big)\alpha_i \right) \ge 1 - \frac{1}{8\ell}
\]

\medskip
by applying quantum phase estimation to 
\[
G_i = \big(2|\psi_i\>\<\psi_i)-I\big) \big( 2\Gamma - I \big) \quad\mbox{and}\quad
|\psi_i\> = |\pi_i\>|0\>
\]

\medskip
\item to achieve the desired accuracy, we have to apply $G_i$ 
\[
\tilde{O}\left(\frac{\ell}{\epsilon}\right)
\]
\end{itemize}
\end{slide}

%%%%%%%%%%%%%%%%%%%%%%%%%%%%%%%%%%%%%%%%%%%%%%%%%%%%%%%%%%%%%%%%%%%%%%%%%

\begin{slide}{Quantum Estimation of the Ratios IV}
\begin{itemize}
\item consider the random variable $Q_0 Q_1 \cdots Q_{\ell-1}$

\item $\Rightarrow$ we have
\[
\Pr\Big( (1-\epsilon)\alpha \le Q \le (1+\epsilon)\alpha \Big) \ge \frac{7}{8} 
\]

\item the success probability decreases to $\ge 3/4$ due to imperfections

\item to obtain the desired accuracy, it suffices to apply operators in $\{W(P_k) : k=0,\ldots,\ell-1\}$
\[
O\Big(
\frac{\ell^2}{\sqrt{\delta}\epsilon}
\Big)
\]
times
\end{itemize}
\end{slide}

%%%%%%%%%%%%%%%%%%%%%%%%%%%%%%%%%%%%%%%%%%%%%%%%%%%%%%%%%%%%%%%%%%%%%%%%%

\begin{slide}{Summary: Quantum Speed-Up}
\begin{itemize}
\item classical complexity
\[
\tilde{O}\Big(
\ell^2 /\big(\delta \epsilon^2\big) 
\Big)
\]

\item quantum complexity
\[
\tilde{O}\Big(
\ell^2/\big(\sqrt{\delta} \epsilon\big)
\Big)
\]

\item $1/\delta \rightarrow 1\sqrt{\delta}$ is due to the quadratic relation between the spectral gaps of $P_i$ and the phase gaps of the corresponding quantum walks $W(P_i)$ 

\medskip
\item $\epsilon^2 \rightarrow \epsilon$ is due to quantum estimation of expected values of observables
\end{itemize}
\end{slide}

%%%%%%%%%%%%%%%%%%%%%%%%%%%%%%%%%%%%%%%%%%%%%%%%%%%%%%%%%%%%%%%%%%%%%%%%%

\begin{slide}{Future Research}

\begin{itemize}
\item improving upon Chebyshev sampling on a quantum computer?

\item quantum speed-up of 

\medskip
\begin{itemize}
\item estimating permanents of matrices with non-negative entries?

\medskip
\item quantum speed-up of estimating the volume of a convex polytope?

\medskip 
\item quantum speed-up of other classical approximation algorithms?
\end{itemize}

\medskip
\item estimating partition functions of \textcolor{red}{quantum} Hamiltonians
\end{itemize}

\end{slide}

%%%%%%%%%%%%%%%%%%%%%%%%%%%%%%%%%%%%%%%%%%%%%%%%%%%%%%%%%%%%%%%%%%%%%%%%%

\end{document}